% Source: https://artofproblemsolving.com/community/c3h1600482_infinity__0
% Stable identity: aops-topic-1600482-q01
\textbf{Error 1: sharing a property does not make two things equal.} Even if both expressions were ``undefined'', that would not imply $0=\infty$. The argument has the form ``$a$ has property $Q$ and $b$ has property $Q$, so $a=b$'', which is invalid: for example $\dfrac{x}{0}$ and $\sqrt{-1}$ (in the real numbers) are both undefined, but nobody concludes that $0=-1$. ``Undefined'' is not a number that both expressions are equal to; it means that no value is assigned at all.

\textbf{Error 2: $\infty$ is not a real number.} The symbol $\infty$ is not a number one can divide by in ordinary arithmetic, so ``$\dfrac{x}{\infty}$'' is not an expression of real-number arithmetic. What is true is a statement about \emph{limits}: $\lim_{t\to\infty}\dfrac{x}{t}=0$ for every fixed real $x$, while $\dfrac{x}{t}$ grows without bound in size as $t\to0$ (for $x\ne0$). So in the sense in which ``$\dfrac{x}{\infty}$'' means anything, it means $0$, not ``undefined'', and the two cases behave in opposite ways.

Hence the argument proves nothing, and $\infty=0$ is not a meaningful (let alone true) statement about real numbers.
